\section{AegisKV Design}
\label{sec:design}

AegisKV separates complete recovery coverage from fast recovery access.
Shared DRAM and shared storage retain complete KV checkpoints; Peer HBM
holds selected hot blocks on a preassigned takeover DP. Its objective is
to reduce workflow completion time (WCT) under finite memory and I/O
capacity while preserving normal cache reuse.

\paragraph{Scope and assumptions.}
We consider fail-stop compute-worker failures with a compatible surviving
DP. The directory, ordered token log and advertised backing copies must
survive the protected failure domain; same-host DRAM or SSD does not
protect against loss of that host. The orchestrator exposes known
fork/join dependencies and accepted inputs. AegisKV preserves the engine's
native attention selections and replay semantics; it does not provide
exactly-once external tool side effects or protection against loss of all
advertised backing domains.

\subsection{Design Insights}
\label{sec:design-insights}

\paragraph{Value recovery in the workflow, without disabling normal reuse.}
A restore advances completion only when it shortens a required path
(Figure~\ref{fig:motivation-recovery-exposure}). However, a high-slack
agent may reuse its history immediately after a tool call. Assigning it
permanently to recomputation would consume shared serving capacity.
AegisKV therefore values both normal reuse and failure recovery, and
separates retaining any complete copy from paying for a faster copy.

\paragraph{Use sparse demand for acceleration, not completeness.}
Native sparse attention accesses only part of historical KV, but later
queries can select different blocks. In Figure~\ref{fig:design-insights},
a causal frequency policy at a 40\% block budget covers 77.8\%, 53.6\%
and 43.0\% of distinct future demand for NOSA, MiniCPM4 and Seer,
respectively. The matched oracle covers 93.7\%, 59.0\% and 43.3\%:
low absolute coverage can reflect limited capacity as well as prediction
error. These observations motivate selective Peer copies backed by a
complete checkpoint and an on-demand miss path.

\begin{figure}[t]
 \centering
 \includegraphics[width=\columnwidth]{figures/retry-kv-prediction.pdf}
 \caption{Future native KV coverage at a 40\% block budget, including
 mandatory blocks. Three PG19 captures per model use queries
 32,257--32,768 to predict 32,769--32,832; an object is one layer's
 64-token block across KV heads. Bars/whiskers are document means/ranges;
 diamonds are same-budget offline oracles. Seer uses the experimental
 Qwen3-8B gate~\cite{song2026seerqwen3}, threshold $10^{-3}$, without
 final-block dense override. Coverage is not recovery speedup.}
 \label{fig:design-insights}
 \Description{Native future-demand coverage varies across models and
 remains below complete recovery coverage.}
\end{figure}

\paragraph{Pay only for transfer that remains exposed.}
Complete DRAM avoids historical reconstruction but still requires H2D.
A Peer copy can remove that transfer; later-layer transfers may already
be hidden by computation that must still execute. Useful residency
therefore depends on each request's committed history, uncommitted tail,
demand reveal time and recovery queues. The Master chooses per-request
layers and hot blocks; the Adapter controls only the total Peer budget.

\subsection{System Overview}
\label{sec:system-overview}

Figure~\ref{fig:system-overview} divides responsibilities across
Application, Controller and Engine. The Application supplies requests,
known dependencies and progress. The Controller contains
\textbf{(1)} Workflow Tracker, \textbf{(2)} Recovery-aware Router,
\textbf{(3)} KV Master and surviving metadata. Tracker supplies duration
estimates and slack; Router binds execution and takeover DPs; Master
chooses complete-checkpoint placement and Peer hot replicas. Within each
DP, \textbf{(4)} Adapter exposes a Peer budget and \textbf{(5)} Connector
executes transfers and recovery. Ordinary execution state, Peer copies
and transient workspace share physical HBM but have separate accounting.

\begin{figure}[t]
 \centering
 \includegraphics[width=\columnwidth]{figures/system-overview.pdf}
 \caption{AegisKV architecture. Black callouts match components (1)--(5)
 in the text; DP1 has the same Adapter/Connector as DP0. The example
 binds DP0 execution to DP1 takeover/Peer placement. Controller state
 includes manifests, token logs and leases. Blue links carry KV,
 dark dashes carry control, red dashes mark retry; orange local blocks
 illustrate an unsealed tail. Complete shared backing lies outside
 both DP failure boundaries. D/S denote DRAM/storage placement.}
 \label{fig:system-overview}
 \Description{Compact Application, Controller and Engine layers with
 numbered components, symmetric DP Local/Peer HBM and shared backing.}
\end{figure}

\paragraph{Objects and checkpoints.}
An immutable object identifies a sealed token block in one layer and TP
shard, together with model/weight version, layout, format and position
semantics. A legal group includes required selector state. Physical
blocks, sparse selection units and kernel tiles need not coincide; the
64-token unit above is a measurement convention, not a universal engine
requirement. Shared physical replicas consume capacity once.

A checkpoint manifest is complete only when every historical layer,
shard, auxiliary state and matching log prefix through frontier $M$ is
committed. A DRAM-resident manifest resolves all dependencies in DRAM;
it is not a hot subset split across tiers. Storage manifests are likewise
complete. Unchanged objects can be shared between checkpoints and only
new objects need copying, but partial extension cannot advance $M$.

Prefill offers newly sealed state; decode triggers review every configured,
block-aligned $Z$ tokens. Arrival, tool handoff, material demand/slack
changes and pressure also trigger review. The acknowledged tail $N$ can
exceed $Z$ when copying lags. Peer coverage is advertised only for objects
with complete committed backing in DRAM or storage.

\subsection{Workflow Tracker: Remaining Time and Slack}
\label{sec:workflow-protection}

Tracker maintains a residual DAG of unfinished execution segments.
Fork/join boundaries separate parent execution from waiting for children,
avoiding double counting. A GBDT estimates remaining duration from causal
features: role/model and current input length; elapsed execution and
completed-call count; and completed calls' output lengths, API times,
context growth and gaps. Missing history is masked. Future tool outputs,
future call counts and the current request's final output length are
excluded. An estimate of remaining time is not reduced by elapsed time
again. Task/tool semantics can be used only when exposed by the application.

Critical-path analysis~\cite{kelley1959cpm} computes earliest and latest
finishes in $O(|V|+|E|)$; their nonnegative difference is segment slack
$\sigma_r$. Requests on one segment share that margin. New requests,
resolved dependencies and execution progress refresh predictions and
publish a new DAG version, invalidating stale placement estimates.

\paragraph{Prediction uncertainty.}
Figure~\ref{fig:tracker-correlations} tests observable inputs and existing
held-out predictions on 1,443 request-2 agents in 170 CC/Weka workflows
\cite{semianalysis2026cctraces}. The target is whole-agent remaining time,
including later calls and gaps. Six feature associations are weak, with
workflow-bootstrap intervals crossing zero. The GBDT's MAE is 406.7\,s
versus 438.6\,s for a training-fold median; the paired difference interval
$[-80.9,4.9]$\,s does not establish a reliable mean improvement. History
features do not consistently help. These results support an explicit
uncertainty guard, not a claim of accurate segment-slack prediction.

\begin{figure}[t]
 \centering
 \includegraphics[width=\columnwidth]{figures/tracker-correlations.pdf}
 \caption{Prediction diagnostics, one sample per agent at request 2.
 (a) Spearman correlations with whole-agent remaining seconds; intervals
 resample 170 whole workflows 2,000 times. Asterisks mark input-length
 proxies. (b) Absolute-error CDFs from five workflow-disjoint folds:
 median baseline, basic GBDT and GBDT with history; all tails are retained
 and the x-axis is logarithmic above 1\,s. This cohort does not validate
 residual-segment Slack.}
 \label{fig:tracker-correlations}
 \Description{Weak feature associations and substantial held-out errors
 motivate conservative handling of duration predictions.}
\end{figure}

When dependencies, segment calibration or profile support are unavailable,
Tracker supplies no positive slack credit. Master falls back to normal
reuse and request recovery cost, without claiming a predicted WCT benefit.
With valid estimates, joint DAG evaluation accounts for shared margins
rather than summing isolated slack savings. Prediction sensitivity and
false-positive slack must be evaluated separately from correlation;
trajectory predictability~\cite{zhang2026heddle} alone does not establish
wall-clock accuracy.

\subsection{Recovery-aware Router}
\label{sec:recovery-router}

Normal dispatch balances load with compatible prefix locality.
Figure~\ref{fig:agent-prefix-reuse} shows why Agent affinity is useful:
97.90\% potential same-agent token coverage and 0.37\% additional
cross-agent coverage in the trace. Cross-agent prefixes can still help
a subagent's first call, so affinity is a preference, not a restriction.
A tool handoff releases execution references and the request's logical
Peer pin; it does not delete reusable backing history.

\begin{figure}[t]
 \centering
 \includegraphics[width=\columnwidth]{figures/agent-prefix-reuse.pdf}
 \caption{Potential prefix reuse over 98,827 requests, 2,090 agents and
 393 CC/Weka workflows. Coverage uses earlier completed requests with
 the same recorded model and workflow. Same/other-agent coverage
 overlaps; only additional cross-agent coverage extends same-agent
 reuse. These are token-prefix opportunities, not runtime cache hits.}
 \label{fig:agent-prefix-reuse}
 \Description{Same-agent prefixes dominate aggregate reuse; cross-agent
 history can also help first subagent calls.}
\end{figure}

At arrival, Router records a versioned, layout-compatible takeover DP
outside the source failure domain. Peer copies reside on that target's
corresponding TP shards. On confirmed failure, Router fences the old
epoch, resolves a complete manifest and activates the bound target subject
to admission and workspace constraints. If it is unavailable, a new
binding can immediately use complete backing with zero Peer coverage;
bytes on the old target are not local hits on the new one.

Before output acknowledgement, a surviving ordered log records request
ID, token position/value and source epoch; prompt and accepted tool inputs
remain recoverable. Replay suppresses duplicate output and rejects stale
callbacks. Rebinding preserves the prior authoritative path until the new
binding is valid. Router queues and target load enter placement cost.

\subsection{KV Master: Three-Tier Placement and Direct Revaluation}
\label{sec:kv-master}

\subsubsection{Full-checkpoint placement and Peer admission}
\label{sec:placement-value}

The Master distinguishes complete-checkpoint DRAM retention, optional
Peer acceleration, and last-copy deletion. New sealed state proposes an
extension in DRAM; if it cannot fit, storage is considered. Existing
immutable dependencies are reused. If neither tier can receive a complete
extension, $M$ stays unchanged and the replay tail grows. A large slack
does not disable caching between normal calls.

For placement $P$ and a feasible candidate $P'$, Master evaluates
\begin{equation}
 \Delta J=A(P',P)+\sum_e p_e\sum_w
       [T_w(P',e)-T_w(P,e)].
 \label{eq:value}
\end{equation}
Here $e$ is a declared continuation scenario over a fixed decision horizon:
normal reuse with no failure, or reuse plus a named source-domain failure.
Scenario probabilities $p_e$ sum to one; raw failure rates are not used as
probabilities. Profiles map committed coverage, tail replay, demand and
shared queues to segment delays; a max-plus DAG pass computes $T_w$.
The transition term $A$ charges foreground copying and temporary occupancy
costs not already included in those scenarios. All terms are in seconds.
If Tracker invalidates the timing model, the fallback replaces the WCT
term with normal-use and request-recovery delays and grants no slack
benefit; it is a degraded policy, not a WCT prediction guarantee.

A common failure includes every interrupted request in that domain.
If two necessary branches are each 10\,s late, accelerating just one to
2\,s still leaves the join 10\,s late. Thus block scores cannot simply be
added. DRAM demotion is compared with actual storage-backed continuation;
Peer is compared with its complete lower path, including exposed H2D;
last-copy deletion is compared with an older checkpoint or reconstruction.
A backing demotion also changes the miss costs of dependent Peer plans.

\begin{figure}[t]
 \centering
 \includegraphics[width=\columnwidth]{figures/retention-value.pdf}
 \caption{Slack alone does not rank retention value. In this isolated-fault
 example, A uses two capacity units to reduce recovery from 4 to 2\,s;
 B uses one to reduce 18 to 8\,s. A at slack 0 saves 1\,s/unit, whereas
 B at slack 6 saves 10\,s/unit. Normal-use value is zero. These constructed
 profiles motivate cost/size comparison, not a policy-speedup claim.}
 \label{fig:retention-value}
 \Description{A higher-slack object can save more workflow time per byte.}
\end{figure}

\subsubsection{Choosing per-request layers and hot blocks}
\label{sec:sparse-placement}

For request $r$, $L_r$ is its Peer prefix depth and $K_{r,l}$ the number
of engine token blocks replicated in layer $l\le L_r$, including required
TP objects. The layer's hot set contains native mandatory dependencies
and blocks ranked by causal access history. Complete backing includes
the remaining blocks and selector state
(Figure~\ref{fig:sparse-recovery}). A prediction miss changes fetching,
not the engine's selection result.

\begin{figure}[t]
 \centering
 \includegraphics[width=\columnwidth]{figures/sparse-recovery.pdf}
 \caption{Complete backing, per-layer Peer budgets and demand misses.
 (a) The DRAM manifest covers B0--B3 in every layer; the tail is uncommitted.
 (b) On the bound target, $L=2$, with $K_1=2$ and $K_2=3$; layer 3 has
 no Peer copy. (c) Layer 1 needs B0/B2: B0 hits Peer and B2 is fetched
 from backing before computation. B3 was unused. The complete manifest
 also covers TP shards and selector state. If demoted as a whole,
 shared storage supplies the miss path before H2D.}
 \label{fig:sparse-recovery}
 \Description{Three complete historical layers, an uncommitted tail,
 different hot-block counts and an explicitly routed layer-1 miss.}
\end{figure}

Master enumerates configured finite grids of layer depths and per-layer
hot-set budgets, including zero Peer and the current placement. It also
allows small bundles across jointly exposed branches. Each candidate's
cost uses current $M$ and a declared distribution or bound for unknown
$N$, including commit lag beyond $Z$. Transfer scheduling accounts for
finite workspace, reveal time and actual remaining computation; avoided
history prefill is not available overlap time. Candidate and victim work
are bounded together. This search does not assume additive values or
claim a globally optimal placement.

Figure~\ref{fig:design-pipeline} illustrates why layer order, prediction
and Peer residency complement each other. The matched comparison is
(c) versus (d): both retain complete DRAM backing and replay the same
64-token tail. Peer saves 24\,ms of exposed waiting in this model; it
neither removes miss fetches nor avoids the required replay computation.
The first seven layers reach the model's demand-only wait floor at this
fixed profile, not a universal depth or a bound under contention.

\begin{figure}[t]
 \centering
 \includegraphics[width=\columnwidth]{figures/recovery-pipeline.pdf}
 \caption{Analytical layered recovery with and without Peer.
 MiniCPM4-8B has 32 layers, 32K history and a
 64-token tail. Replay is calibrated at 2.870\,ms/layer; assumed effective
 H2D bandwidth is 6.25\,GB/s (5.369\,ms per full layer). Trace-derived
 mean preload and preload-plus-miss fractions are 0.6584/0.7748.
 (a) waits for bulk restoration; (b) overlaps full-layer restoration;
 (c) preloads hot blocks and fetches revealed misses; (d) adds 147.5\,MiB
 of Peer hot sets in the first seven layers. Numbers are tail completion
 and exposed wait. Cold copies follow replay; the ideal preemptible
 stream and homogeneous layer costs omit bounded-workspace and queue
 effects. This is not end-to-end fault-recovery measurement.}
 \label{fig:design-pipeline}
 \Description{The same required replay computation with progressively
 reduced exposed KV waits: 172, 83, 44 and 20 milliseconds.}
\end{figure}

Neither $L$ nor $K$ must increase on every decode step. A fixed sparse
budget can preserve $K$ while changing block identities; a proportional
budget can increase demand, but Peer admission remains capacity-limited.
At each $Z$ review or material event, the Master can retain or replace the
current set. Proposed and committed sets remain distinct until transfers
finish. If no exposed wait is removed, $L_r=0$ is valid even at zero slack.

\subsubsection{Online admission, direct revaluation, and eviction}
\label{sec:placement-plan}

Each decision snapshots DAG and binding versions, causal demand, complete
manifests, queues, leases and physical byte ledgers. DRAM/storage candidates
are complete manifests; Peer candidates are legal layer/block bundles.
The safe victim set excludes physical readers and copy dependencies.
Normal Peer competition also respects running requests' logical pins;
request completion or tool handoff unpins them. Self hot-set replacement
and Adapter pressure may override logical pins, never physical leases.

For each candidate, a bounded rotating scan examines eligible victims
in stable-ID order. It appends the feasible removal with least current
whole-plan loss per physical byte released, breaking ties by fewer moved
bytes and then ID. The tentative placement is revalued after every
removal. A manifest releasing no bytes alone is skipped unless a legal
owner bundle can release its shared objects. A DRAM victim needing
storage space includes lower-tier relief and fallback changes in the
same plan. Search stops when space exists for the candidate and its
workspace, or the work budget is exhausted.

The ratio constructs a victim set; the final decision evaluates the
complete candidate-plus-victim plan with Equation~\ref{eq:value}.
Optional placement must improve cost beyond uncertainty and migration
margins. These are configured guards, not statistical confidence bounds.
Forced capacity relief can accept positive cost and has no incoming
candidate. If no safe plan fits, admission pauses; an infeasible plan
never changes current coverage or makes reserved bytes free.

\par\medskip\noindent\begin{minipage}{\columnwidth}
\raggedright
 \refstepcounter{designalgorithm}\label{alg:placement}
 \textbf{Algorithm \thedesignalgorithm: Direct placement revaluation}
 \par\smallskip
 \begin{tabularx}{\columnwidth}{@{}r@{\quad}X@{}}
 \toprule
 1 & Snapshot DAG/binding versions, manifests, demand, queues and capacity \\
 2 & Refresh affected costs; use conservative fallback if profiles are invalid \\
 3 & On required shrink, build a safe current-loss relief plan \\
 4 & Otherwise enumerate bounded full-checkpoint and Peer L/K candidates \\
 5 & For each candidate, revalue victims until space fits or work is exhausted \\
 6 & Include lower-tier copies, victims, fallback changes and workspace \\
 7 & Compare complete plans; keep the best improvement beyond margins \\
 8 & If no plan qualifies, retain state; pause admission under pressure \\
 9 & Recheck versions and reserve immediately executable capacity \\
 10 & Enqueue copy--commit--redirect--drain--release dependencies \\
 11 & Reconcile completion or failure; trigger a fresh review \\
 \bottomrule
 \end{tabularx}
\end{minipage}\par\medskip

Cached estimates are reused only while their dependencies remain valid.
Profiling, victim comparisons and DAG passes all count against the work
budget. Unsupported or stale resource profiles freeze optional promotions;
necessary relief uses bounded recency order under the same safety rules.
Direct revaluation has no rent, credit or renewal state; the rejected
rental alternative is evaluated separately in Section~\ref{sec:eval-rental}.

\subsection{Adapter: Windowed Peer Capacity}
\label{sec:adapter}

Each DP's Adapter sets a Peer budget for the next time window from the
managed HBM pool, expected serving allocation and reserved headroom,
capped by a configured maximum and clipped at zero. Serving allocation
includes execution state and near-term demand, without subtracting Peer
occupancy twice. The window duration and cap are configurable; $Z$ is
fixed independently. Adapter does not set per-request $L$ or $K$.

A lower target pauses new Peer admission while actual plus reserved
occupancy exceeds the target. Master selects a safe relief plan, including
logically pinned replicas when necessary. A valid immutable backing copy
allows clean Peer release without redundant D2H. Active reads or copies
must drain before memory becomes available. Emergency foreground pressure
can request relief within a window; if leases prevent immediate release,
serving applies backpressure rather than reusing occupied addresses.

\subsection{Connector: Committed Copies and Layered Recovery}
\label{sec:connector}

Connector applies reservations and bounded transfer queues for foreground
reuse, demanded recovery and checkpoint movement. Configured service
shares and transfer quanta limit interference; real misses precede
speculative reads within the recovery share. Copying and Peer occupancy
remain costs in the placement objective.

\paragraph{Publish complete state before releasing its predecessor.}
A full extension commits only after validation of every required layer,
TP shard, selector object and corresponding token-log prefix. Moving a
complete DRAM checkpoint to storage follows copy, backend persistence,
manifest commit, reader redirection, reader drainage and source release.
Partial writes cannot release the source; duplicate or stale-epoch
callbacks cannot advance coverage. Temporary duplicate objects and
workspace count against real capacity, and a buffered write acknowledgement
alone is not proof of the backend's persistence contract.

Before last-copy deletion, Connector publishes a complete older manifest
or a prompt-and-token-log reconstruction fallback. Peer claims lacking
complete backing are withdrawn in the same transition. On copy failure,
the prior authoritative location remains valid and failed reservations
are reconciled; pending copies never count as current protection.

\paragraph{Restore according to actual demand.}
Recovery acquires the manifest, binding and location leases. Compatible
Peer blocks on the bound DP bypass H2D; missing blocks use complete DRAM,
or shared-storage reads followed by H2D. Layer-ordered prefetch overlaps
only with remaining replay computation. Actual selection misses block
the affected operation until inputs arrive; dense layers use full demand
and missing selector state is restored before selection. TP collectives
and replay batching preserve native semantics. Adopted Peer objects move
to execution accounting without double-counting or freeing their bytes.
Prefill failures use the last complete frontier, possibly zero.
